% ------------------------------------------------------------------
%  Prefazione (non numerata: \chapter nel frontmatter)
%  Da riscrivere a libro finito. Materiale: ricerca 01 §1, §17.3.
% ------------------------------------------------------------------
\chapter{Prefazione}

\iniziale{Q}{uesto libro} nasce da una domanda semplice: a che cosa
serve studiare, se poi dimentichiamo quasi tutto? Chiunque sia passato
per la scuola conosce la sensazione di aver superato una verifica o un
esame e, poche settimane dopo, di non ricordarne più nulla. Non è una
colpa personale: è il modo in cui funziona la memoria quando la si usa
male.

Le pagine che seguono si rivolgono anzitutto a chi studia
all'università --- dalla matricola al laureando --- e a chi vi insegna:
è all'università che si studia più da soli, ed è lì che si formano gli
insegnanti di domani.
Raccolgono ciò che la ricerca sull'apprendimento ha chiarito negli
ultimi decenni --- e ciò che i filosofi e i maestri antichi avevano
intuito molto prima --- e lo traducono in pratiche quotidiane,
verificabili, alla portata di tutti.

\segnaposto{da completare a libro finito: perché questo libro, perché comincia dall'università (e la scuola verrà in una seconda fase), perché è gratuito, che cosa troverà il lettore}

\finecapitolo
